\paragraph{Determinants before degrees.}
The two filtrations give more than a degree identity. They force a
finite tensor power of the determinant to be trivial, even on a
singular base.

\begin{proposition}\label{geo-frobenius-determinant-torsion}
Let $X$ be a scheme over $k=\mathbf F_{p^f}$, let $F$ be absolute
Frobenius, and let $A$ be a finite commutative $\mathbf F_p$-algebra.
Put $B=k\otimes_{\mathbf F_p}A$. Suppose finite locally free
$\mathcal O_X$-modules with $A$-actions have $A$-linear exact sequences
\begin{equation}\label{geo-hodge-conjugate-sequences}
\begin{gathered}
0\longrightarrow\mathcal U\longrightarrow\mathcal H
 \longrightarrow\mathcal V\longrightarrow0,\\
0\longrightarrow F^*\mathcal V\longrightarrow\mathcal H
 \longrightarrow F^*\mathcal U\longrightarrow0.
\end{gathered}
\end{equation}
The pullbacks carry the pulled-back $A$-actions. For any idempotent
$e\in B$, put $\mathcal D_e=\det(e\mathcal H)$. Then
\[
                    \mathcal D_e^{\otimes(p^f-1)}
                         \simeq\mathcal O_X.
\]
Let $\chi:B\to k$ be a $k$-algebra character and $B_e=eB$ its
Artin local factor, of length $n=\dim_k B_e$. If $e\mathcal H$ is
locally free over $\mathcal O_X\otimes_k B_e$, then
$\mathcal Q_\chi=\mathcal H\otimes_{B,\chi}k$ is a vector bundle and
\begin{equation}\label{geo-character-determinant-torsion}
       (\det\mathcal Q_\chi)^{\otimes n(p^f-1)}
                             \simeq\mathcal O_X.
\end{equation}
In particular every pullback of $\mathcal Q_\chi$ to a proper curve
has degree zero.
\end{proposition}
\begin{proof}
The automorphism $\sigma(c\otimes a)=c^p\otimes a$ of $B$ fixes $A$;
it need not equal Frobenius on the possibly nonreduced algebra $B$.
For any bundle $\mathcal W$ with this action,
$F^*(e\mathcal W)=\sigma(e)F^*\mathcal W$. Project the second
sequence by $\sigma(e)$ and the first by $e$, and take determinants:
\[
 \mathcal D_{\sigma(e)}
 \simeq F^*\{\det(e\mathcal V)\otimes\det(e\mathcal U)\}
 \simeq F^*\mathcal D_e
 \simeq\mathcal D_e^{\otimes p}.
\]
The last identity follows directly by raising transition functions
to their $p$-th powers. Since $\sigma^f=1$, iteration proves the
first assertion. Idempotent images are direct summands, and the
bundle sequences are locally split; no flatness of $F$ is required.

The residue field of $B_e$ is $k$. Tensor a composition series of
$B_e$ with the locally free coefficient module $e\mathcal H$.
Its $n$ quotient bundles are all isomorphic to $\mathcal Q_\chi$,
so $\mathcal D_e\simeq(\det\mathcal Q_\chi)^{\otimes n}$.
This proves \eqref{geo-character-determinant-torsion}.
Degree on a proper curve is additive under tensor products and
takes integer values, giving the final assertion. Rank zero causes
no exception.
\end{proof}

The factor $n$ retains the nilpotent multiplicity at a ramified
coefficient prime. The proposition applies to the entire base only
if its filtrations and coefficient module are defined there. If
they are constructed only on test curves, its conclusion is
curvewise numerical triviality, which already suffices below.

\paragraph{One finite valuation on the whole fibre.}
We use the elementary fact that a numerically trivial line bundle on
a proper integral variety has no nonzero section with a zero.
Indeed Chow's lemma reduces to a projective integral modification.
A nonempty effective zero divisor then meets a projective curve
not contained in it, after a finite field extension if necessary.
The restricted section has positive degree. Pullback preserves
numerical triviality, by the degree formula for maps of proper
curves, giving a contradiction
\cite[Tags 0200 and 0B40]{stacks-finite-place}.

\begin{theorem}\label{geo-vertical-line-rigidity}
Let $R$ be a discrete valuation ring with uniformizer $\pi$,
fraction field $K$ and residue field $k$. Let $X$ be a normal
integral proper flat $R$-scheme with nonempty connected special fibre.
Let $\mathcal L$ be a line bundle whose restriction to every proper
integral curve in $X_k$ has degree zero. Suppose $s$ generates
$\mathcal L$ everywhere on $X_K$. Then there are integers $M>0$
and $a$ for which
\begin{equation}\label{geo-vertical-generator}
                    s^{\otimes M}/\pi^a
          \quad\hbox{generates }\mathcal L^{\otimes M}
          \quad\hbox{on all of }X.
\end{equation}
Write the fibre divisor as $\operatorname{div}(\pi)=\sum_i m_iX_i$.
The normalized order $q=a/M$ is unique, and its denominator in
lowest terms divides $\gcd_i m_i$. In particular $M=1$ is possible
when the fibre multiplicities are relatively prime.
\end{theorem}
\begin{proof}
The divisor $D=\operatorname{div}(s)$ is vertical, so write
$D=\sum_i d_iX_i$ and set $q=\min_i(d_i/m_i)=a/M$ in lowest terms,
with $M>0$. The Cartier divisor
\[
                  E=MD-a\operatorname{div}(\pi)
\]
has nonnegative coefficients and omits at least one component.
Normality implies that it is effective: locally a rational
function with nonnegative order at every prime divisor is regular.

Suppose $E$ is nonempty. The intersection graph of the irreducible
components of the connected special fibre is connected. Hence
some component $X_i$ omitted from $E$ meets its support. Restricting
$E$ to the reduced integral scheme $X_i$ gives a nonempty effective
Cartier divisor. Its line bundle is
$\mathcal L^{\otimes M}|_{X_i}$, because the whole fibre is
principal. This line bundle is numerically trivial, contradicting
the preceding fact. Thus $E=0$. Normality again implies that a
rational section with neither zeros nor poles is a generator,
proving \eqref{geo-vertical-generator}.

All ratios $d_i/m_i$ are therefore $q$. In lowest terms,
$M d_i=a m_i$ and $\gcd(a,M)=1$, so $M$ divides every $m_i$.
These identities also prove uniqueness.
\end{proof}

The theorem does not assert $q=0$. For example, the trivial line
with generic generator $\pi$ has normalized order one. A totally
ramified extension of degree $m$, with a uniformizer as generator,
can give order $1/m$, so the multiplicities cannot be discarded.
Connectedness is also essential: the generators $1$ and $\pi$
on two disjoint copies of $\operatorname{Spec}R$ have different
orders.

For the prescribed conjugate determinant, the characteristic-$p$
proposition supplies the curve degrees once the actual Hodge
filtrations and coefficient freeness are identified. The theorem
then makes the finite norm constant across the connected fibre.
Its value still requires the same integral complex multiplication
evaluation and different correction as in
\eqref{geo-exact-conjugate-degree}. Neither that evaluation nor
the boundary lattice identification follows from the abstract
determinant argument.
